\documentclass[10pt,a4paper]{amsart}
\usepackage[margin=19mm]{geometry}
\usepackage{amsmath,amssymb,mathtools}
\usepackage[scr=rsfs,cal=euler]{mathalfa}
\usepackage{xfrac}
\usepackage[unicode,hidelinks,bookmarks=false]{hyperref}
\newcommand{\ie}{i.e.\ }
\newcommand{\cf}{cf.\ }
\newcommand{\resp}{respectively\ }
\newcommand{\Subsection}{\subsection}
\providecommand{\nobreakdash}{\nobreak\hbox{-}}
\setlength{\emergencystretch}{2em}
\multlinegap0pt
\usepackage{amssymb}
\usepackage[shortlabels]{enumitem}
\usepackage{float}
\usepackage{bm}
\newtheorem{theorem}{Theorem}
\newtheorem{proposition}{Proposition}
\title{Shortest LCD embeddings of binary, ternary and quaternary linear codes}
\author{Junmin An, Ji-Hoon Hong, Jon-Lark Kim, Haeun Lim}
\date{Source: arXiv:2601.20600v2 (CC BY 4.0)}
\begin{document}
\maketitle

\noindent\emph{Source and licence.} Shortest LCD embeddings of binary, ternary and quaternary linear codes. Version arXiv:2601.20600v2 (CC BY 4.0), distributed by arXiv under the Creative Commons Attribution 4.0 International licence, http://creativecommons.org/licenses/by/4.0/, as recorded in the arXiv licence metadata for this exact version and shown on the abstract page https://arxiv.org/abs/2601.20600v2. Full licence notice: \texttt{licenses/arxiv-2601.20600-CC-BY-4.0.txt}.\par\smallskip

\noindent\textit{Editorial scope.} Counted unit: the article's proposition proposition (source0.tex lines 277--291). The article's complete original proof of it is delivered with it (source0.tex lines 292--320, one \texttt{proof} environment of 29 lines). No mathematics is written, repaired or re-derived: the statement and the proof are the article's own lines, in the article's own notation, and no retained line is substituted or re-flowed.  Retained uncounted as the source-local context the proof needs: lines 238--252.  Nothing was omitted from any retained span, so the package carries no omission record.  No retained line contains a citation, so the wrapper carries no bibliography and no bibliography entry.  No retained line contains a figure environment, an image-inclusion command or drawing code, so no image is delivered and the package declares no asset dependency.  Editorial formatting only, in addition: an AMS \texttt{amsart} preamble that restates the article's own declarations and macros used by the retained lines, namely the environments \texttt{theorem}, \texttt{proposition} on separate counters, together with the standard packages those lines need.  The retained environments print in this document as Theorem 1, Proposition 1 in order of appearance, whereas the article prints the same items with the numbering of its own full document.  The complete native source of the article as distributed in the arXiv e-print for this exact version is bundled unchanged as \texttt{sources/193537/source0.tex}.
\begin{theorem}\label{construct-LCD}
    Let $\mathcal{C}$ be an $[n, k]$ code over $F$ with $\ell=\dim({\rm Hull}(\mathcal{C}))$ and generator matrix
    \[
    G=\begin{bmatrix}
    G({\rm Hull}(\mathcal{C}))\\A
    \end{bmatrix}.
    \]
    Let $\tilde{G}$ be a matrix obtained by appending an invertible $\ell\times \ell$ matrix $D$ and a $(k-\ell)\times \ell$ matrix $C$ as follows:
    \[
    \tilde{G}=\begin{bmatrix}
        G({\rm Hull}(\mathcal{C})) & D\\A & C
    \end{bmatrix}.
    \]
    Then $\tilde{G}$ generates a shortest LCD embedding of $\mathcal{C}$. Conversely, every shortest LCD embedding of $\mathcal{C}$ has a generator matrix of this form.
\end{theorem}

\begin{proposition}
    Let $\mathcal{C}$ be an $[n, k]$ code over $F$ with $\ell=\dim({\rm Hull}(\mathcal{C}))$ and $\tilde{\mathcal{C}}$ be a shortest LCD embedding of $\mathcal{C}$. Then $\tilde{\mathcal{C}}$ has a generator matrix of the form
    \[
    \tilde{G}=\begin{bmatrix}
        G({\rm Hull}(\mathcal{C}))' & I_\ell\\A' & \mathcal{O}
    \end{bmatrix},
    \]
    where $G({\rm Hull}(\mathcal{C}))'$ generates ${\rm Hull}(\mathcal{C})$ and
    \[
    G=\begin{bmatrix}
        G({\rm Hull}(\mathcal{C}))'\\A'
    \end{bmatrix}
    \]
    is a generator matrix of $\mathcal{C}$.
\end{proposition}

\begin{proof}
    By Theorem~\ref{construct-LCD}, $\tilde{\mathcal{C}}$ has a generator matrix of the form
    \[
    \tilde{G}_0=\begin{bmatrix}
        G(\text{Hull}(\mathcal{C})) & D\\A & C
    \end{bmatrix},
    \]
    where $G(\text{Hull}(\mathcal{C}))$ is a generator matrix of $\text{Hull}(\mathcal{C})$ and $D$ is an invertible matrix. Define
     \[
    M=\begin{bmatrix}
        I_\ell & \mathcal{O}\\-C & I_{k-\ell}
    \end{bmatrix}\quad\mbox{and}\quad N=\begin{bmatrix}
        D^{-1} & \mathcal{O}\\\mathcal{O} & I_{k-\ell}
    \end{bmatrix}.
    \]
    Since both $M$ and $N$ are invertible,
    \[
    \tilde{G}=MN\tilde{G}_0=\begin{bmatrix}
        D^{-1}G(\text{Hull}(\mathcal{C})) & I_\ell\\A-CD^{-1}G(\text{Hull}(\mathcal{C})) & \mathcal{O}
    \end{bmatrix}
    \]
    generates $\tilde{\mathcal{C}}$. It follows that the matrix
    \[
    G=\begin{bmatrix}
        D^{-1}G(\text{Hull}(\mathcal{C}))\\A-CD^{-1}G(\text{Hull}(\mathcal{C}))
    \end{bmatrix}
    \]
    which is a punctured matrix of $\tilde{G}$ on the last $\ell$ columns generates $\mathcal{C}$ such that $D^{-1}G(\text{Hull}(\mathcal{C}))$ generates $\text{Hull}(\mathcal{C})$.
\end{proof}

\end{document}
